\section{Convolutional Neural Networks in R1CS}

\subsection{Overview}

Convolutional layers are structured linear operators with parameter sharing.
From the perspective of zero-knowledge proofs, they are compiled into constraints
in exactly the same way as fully connected layers, but with a specific indexing structure.

\medskip

\noindent
\textbf{Key idea.}
A convolution is a linear map:
\[
y = \text{Conv}(x; K) + b,
\]
so each output entry becomes an affine constraint in R1CS.

\subsection{A Minimal CNN Example}

Consider a single-channel $3 \times 3$ input:
\[
X =
\begin{bmatrix}
x_{11} & x_{12} & x_{13} \\
x_{21} & x_{22} & x_{23} \\
x_{31} & x_{32} & x_{33}
\end{bmatrix}.
\]

Let the convolution kernel be $2 \times 2$:
\[
K =
\begin{bmatrix}
k_{11} & k_{12} \\
k_{21} & k_{22}
\end{bmatrix},
\quad \text{bias } b.
\]

We use stride $1$ and no padding. The output is $2 \times 2$:
\[
Y =
\begin{bmatrix}
y_{11} & y_{12} \\
y_{21} & y_{22}
\end{bmatrix}.
\]

\subsection{Convolution as Linear Computation}

Each output entry is:
\[
y_{11} = k_{11}x_{11} + k_{12}x_{12} + k_{21}x_{21} + k_{22}x_{22} + b,
\]
\[
y_{12} = k_{11}x_{12} + k_{12}x_{13} + k_{21}x_{22} + k_{22}x_{23} + b,
\]
\[
y_{21} = k_{11}x_{21} + k_{12}x_{22} + k_{21}x_{31} + k_{22}x_{32} + b,
\]
\[
y_{22} = k_{11}x_{22} + k_{12}x_{23} + k_{21}x_{32} + k_{22}x_{33} + b.
\]

\subsection{Witness Vector}

Flatten all variables into a witness:
\[
w = (1,\; x_{11},x_{12},\ldots,x_{33},\; y_{11},y_{12},y_{21},y_{22}).
\]

\subsection{R1CS Constraints for Convolution}

Each output pixel gives one affine constraint:
\[
(\text{linear combination of inputs} + b)\cdot 1 = y_{ij}.
\]

For example, for $y_{11}$:
\[
(k_{11}x_{11} + k_{12}x_{12} + k_{21}x_{21} + k_{22}x_{22} + b)\cdot 1 = y_{11}.
\]

Similarly for the other entries.

\medskip

\noindent
\textbf{Observation.}
Convolution does not introduce new types of constraints.  
It is simply a sparse, structured linear map.

\subsection{Adding Nonlinearity}

Suppose we apply a polynomial activation:
\[
h_{ij} = y_{ij}^2.
\]

Each activation becomes:
\[
y_{ij} \cdot y_{ij} = h_{ij}.
\]

Thus we add one multiplicative constraint per output entry.

\subsection{Full Constraint Set}

For each output location $(i,j)$:
\begin{itemize}
\item One affine constraint:
\[
(\text{weighted patch} + b)\cdot 1 = y_{ij}
\]
\item One multiplication constraint:
\[
y_{ij} \cdot y_{ij} = h_{ij}
\]
\end{itemize}

\subsection{Matrix View (Im2Col Perspective)}

We can rewrite convolution as matrix multiplication.

Define:
\[
\text{vec}(K) = (k_{11}, k_{12}, k_{21}, k_{22}),
\]

and extract patches:
\[
p_1 = (x_{11},x_{12},x_{21},x_{22}),
\]
\[
p_2 = (x_{12},x_{13},x_{22},x_{23}),
\]
\[
p_3 = (x_{21},x_{22},x_{31},x_{32}),
\]
\[
p_4 = (x_{22},x_{23},x_{32},x_{33}).
\]

Then:
\[
y_i = \langle \text{vec}(K), p_i \rangle + b.
\]

Thus convolution reduces to:
\[
y = W_{\text{conv}} x + b,
\]
where $W_{\text{conv}}$ is a sparse matrix with repeated parameters.

\subsection{Generalization}

For a general convolution layer:

\begin{itemize}
\item Input: $C_{\text{in}} \times H \times W$
\item Kernel: $C_{\text{out}} \times C_{\text{in}} \times k \times k$
\item Output: $C_{\text{out}} \times H' \times W'$
\end{itemize}

Each output element is:
\[
y_{c,i,j} = \sum_{c'} \sum_{u,v}
K_{c,c',u,v} \cdot x_{c',i+u,j+v} + b_c.
\]

Each such equation becomes one affine R1CS constraint.

\subsection{Complexity}

Let $N = C_{\text{out}} H' W'$ be the number of output elements.

\begin{itemize}
\item Affine constraints: $N$
\item Multiplication constraints (activation): $N$
\item Total constraints: $O(N \cdot k^2 \cdot C_{\text{in}})$ operations
\end{itemize}

\subsection{Remarks}

\begin{itemize}
\item Convolution is efficient in R1CS due to parameter sharing.
\item The main cost arises from the number of output elements.
\item Nonlinearities dominate constraint complexity in practice.
\end{itemize}

\medskip

\noindent
\textbf{Key Insight.}  
A convolutional network is still a composition of linear maps and pointwise nonlinearities.  
Thus, it naturally compiles into R1CS with no new primitives beyond those used for MLPs.